% Zbiór zadań — rozdział 2
% Numer rozdziału wynika z nazwy tego pliku.

\begin{exo}[Od zimnego lodu do wrzącej wody: rzędy wielkości]{chauffage-glace-eau-ebullition}
\seoready{true}
\keywords{kalorymetria, przemiana fazowa, topnienie, parowanie, ciepło utajone, rząd wielkości}

\textbf{Dane liczbowe (przy ciśnieniu atmosferycznym):}
\[
\begin{aligned}
c_{\text{lodu}} &= 2{,}1\times10^3\,\text{J}\!\cdot\!\text{kg}^{-1}\!\cdot\!\text{K}^{-1},
& L_f &= 3{,}34\times10^5\,\text{J}\!\cdot\!\text{kg}^{-1},\\
c_{\text{wody}} &= 4{,}18\times10^3\,\text{J}\!\cdot\!\text{kg}^{-1}\!\cdot\!\text{K}^{-1},
& L_v &= 2{,}26\times10^6\,\text{J}\!\cdot\!\text{kg}^{-1}.
\end{aligned}
\]
Przy ciśnieniu atmosferycznym stopniowo ogrzewa się kilogram lodu o temperaturze początkowej $-100\,{}^\circ\text{C}$.

\begin{enumerate}
  \item Obliczyć energię potrzebną do ogrzania lodu od $-100\,{}^\circ\text{C}$ do $0\,{}^\circ\text{C}$.
  \item Obliczyć energię potrzebną do całkowitego stopienia lodu w temperaturze $0\,{}^\circ\text{C}$.
  \item Obliczyć energię potrzebną do ogrzania ciekłej wody od $0\,{}^\circ\text{C}$ do $100\,{}^\circ\text{C}$.
  \item Obliczyć dodatkową energię potrzebną do całkowitego odparowania wody w temperaturze $100\,{}^\circ\text{C}$.
  \item Który etap wymaga najwięcej energii? Porównać energię potrzebną do ogrzania kilograma wody od $0\,{}^\circ\text{C}$ do $100\,{}^\circ\text{C}$ z energią potencjalną masy $m$ spadającej z wysokości dziesięciu metrów. Jaka masa $m$ musi spaść, aby uwolnić taką samą energię?
  \item Rozważyć teraz proces odwrotny. Masa $m_v$ pary wodnej skrapla się na szybie samochodowej, a otrzymana woda nie ochładza się następnie. Wyrazić ciepło $Q_{\text{wyd}}$ wydzielone podczas skraplania i obliczyć je dla $m_v=1{,}0\times10^{-2}\,\text{kg}$. Dla temperatury szyby przyjąć $L_v=2{,}45\times10^6\,\text{J}\!\cdot\!\text{kg}^{-1}$.
\end{enumerate}

\begin{indication}
Do ogrzewania fazy bez przemiany fazowej użyć $Q=mc\Delta T$. Dla przemiany fazowej w stałej temperaturze użyć $Q=mL$. Energia potencjalna grawitacyjna masy $m$ na wysokości $h$ wynosi $E_p=mgh$; przyjąć $g=9{,}81\,\text{m}\!\cdot\!\text{s}^{-2}$.
\end{indication}

\begin{solution}
\textbf{1.} Ogrzanie lodu do temperatury topnienia wymaga
\[
Q_1=mc_{\text{lodu}}\Delta T
=1{,}0\times2{,}1\times10^3\times100
=2{,}1\times10^5\,\text{J}.
\]

\textbf{2.} Podczas topnienia temperatura pozostaje równa $0\,{}^\circ\text{C}$:
\[
Q_2=mL_f=1{,}0\times3{,}34\times10^5
=3{,}34\times10^5\,\text{J}.
\]

\textbf{3.} Do ogrzania ciekłej wody od $0$ do $100\,{}^\circ\text{C}$ potrzeba
\[
Q_3=mc_{\text{wody}}\Delta T
=1{,}0\times4{,}18\times10^3\times100
=4{,}18\times10^5\,\text{J}.
\]

\textbf{4.} Całkowite odparowanie wymaga ponadto
\[
Q_4=mL_v=1{,}0\times2{,}26\times10^6
=2{,}26\times10^6\,\text{J}.
\]
Sam ten etap wymaga więc energii rzędu kilku megadżuli.

\textbf{5.} Parowanie jest zdecydowanie najbardziej energochłonnym etapem:
\[
Q_4=2{,}26\times10^6\,\text{J}
>Q_3=4{,}18\times10^5\,\text{J}
>Q_2=3{,}34\times10^5\,\text{J}
>Q_1=2{,}1\times10^5\,\text{J}.
\]
Parowanie wymaga w szczególności około $(2{,}26\times10^6)/(4{,}18\times10^5)\approx5{,}4$ razy więcej energii niż ogrzanie ciekłej wody od $0$ do $100\,{}^\circ\text{C}$.

Dla masy $m$ spadającej z wysokości $h=10\,\text{m}$ mamy $E_p=mgh$. Z warunku $mgh=Q_3$ wynika
\[
m=\frac{Q_3}{gh}
=\frac{4{,}18\times10^5}{9{,}81\times10}
\approx4{,}26\times10^3\,\text{kg}.
\]
Z wysokości dziesięciu metrów musiałaby zatem spaść masa około $4{,}3$ tony, aby uwolnić tyle energii, ile potrzeba do ogrzania jednego kilograma wody od $0$ do $100\,{}^\circ\text{C}$! To ogromna wartość i wyjaśnia, dlaczego ogrzewanie wody stanowi istotną część domowego zużycia energii. Woda ma także bardzo duże ciepło właściwe w porównaniu z typowymi metalami; jest ono na przykład około dziesięć razy większe niż dla miedzi (zob. następne zadania).

\textbf{6.} Skraplanie jest procesem odwrotnym do parowania: para przekazuje szybie ciepło utajone
\[
Q_{\text{wyd}}=m_vL_v.
\]
Dla $m_v=1{,}0\times10^{-2}\,\text{kg}$,
\[
Q_{\text{wyd}}=1{,}0\times10^{-2}\times2{,}45\times10^6
=2{,}45\times10^4\,\text{J}.
\]
Skroplenie nawet niewielkiej masy pary może więc wydzielić niepomijalną ilość ciepła, które zostaje pochłonięte przez szybę i jej otoczenie.
\end{solution}
\end{exo}

\begin{exo}[Ewapotranspiracja dużego drzewa: naturalny klimatyzator?]{odg-evapotranspiration-arbre}
\seoready{true}
\keywords{rząd wielkości, ewapotranspiracja, drzewo, ciepło utajone, moc chłodnicza, klimatyzator, COP}

W gorący i suchy dzień duże drzewo, mające dostateczny dostęp do wody, może w wyniku ewapotranspiracji oddać około $500\,\text{L}=5{,}0\times10^{-1}\,\text{m}^3$ wody. Ponieważ transpiracja zachodzi głównie przy świetle, przyjmujemy, że trwa przez dwanaście godzin. Korona drzewa pokrywa na ziemi powierzchnię $A_{\text{drzewa}}=50\,\text{m}^2$. Dane są gęstość wody i jej ciepło parowania przy ciśnieniu atmosferycznym:
\[
\rho_{\text{wody}}=1{,}0\times10^3\,\text{kg}\!\cdot\!\text{m}^{-3},
\qquad L_v=2{,}45\times10^6\,\text{J}\!\cdot\!\text{kg}^{-1}.
\]

\begin{enumerate}
  \item Obliczyć masę wody oddanej przez drzewo w ciągu dnia oraz energię potrzebną do jej odparowania.
  \item Wyjaśnić, dlaczego proces ten powoduje chłodzenie.
  \item Obliczyć średnią moc chłodniczą drzewa na metr kwadratowy podczas dwunastu godzin aktywnej transpiracji.
  \item Dla porównania rozważyć klimatyzator pobierający moc elektryczną $P_{\mathrm{el}}=2{,}0\times10^3\,\text{W}$, o współczynniku wydajności chłodniczej $\mathrm{COP}=3{,}0$, chłodzący pomieszczenie o powierzchni $A_{\text{pom}}=20\,\text{m}^2$. Z definicji $\mathrm{COP}=P_{\text{cool}}/P_{\mathrm{el}}$. Obliczyć moc chłodniczą oraz moc chłodniczą na metr kwadratowy i porównać z wynikami dla drzewa.
  \item Dlaczego nie można stwierdzić, że drzewo i klimatyzator dają dokładnie takie samo odczuwalne chłodzenie, nawet jeśli ich moce na jednostkę powierzchni mają ten sam rząd wielkości?
  \item Czy można chłodzić mieszkanie za pomocą wielu roślin doniczkowych? (Wiedza ogólna: odpowiedź zależy od procesów biologicznych.)
\end{enumerate}

\begin{indication}
Użyć $m=\rho V$, $Q_{\mathrm{par}}=mL_v$ i $P=Q/\tau$. Dla klimatyzatora $P_{\text{cool}}=\mathrm{COP}\,P_{\mathrm{el}}$.
\end{indication}

\begin{solution}
\textbf{1.} Odpowiadająca objętości masa wody wynosi
\[
m=\rho_{\text{wody}}V
=1{,}0\times10^3\times5{,}0\times10^{-1}
=5{,}0\times10^2\,\text{kg}.
\]
Energia potrzebna do jej odparowania jest równa
\[
Q_{\mathrm{par}}=mL_v
=5{,}0\times10^2\times2{,}45\times10^6
=1{,}225\times10^9\,\text{J}.
\]
Jest to energia rzędu jednego gigadżula.

\textbf{2.} Parowanie jest endotermiczną przemianą fazową: potrzebna energia jest pobierana z liści, gleby i otaczającego powietrza. Odbieranie ciepła obniża temperaturę liści oraz ich najbliższego otoczenia.

\textbf{3.} Dwanaście godzin odpowiada $\tau=12\times3600=4{,}32\times10^4\,\text{s}$. Średnia moc chłodnicza wynosi zatem
\[
P_{\text{drzewa}}=\frac{Q_{\mathrm{par}}}{\tau}
=\frac{1{,}225\times10^9}{4{,}32\times10^4}
\approx2{,}84\times10^4\,\text{W},
\]
a na jednostkę powierzchni
\[
\frac{P_{\text{drzewa}}}{A_{\text{drzewa}}}
=\frac{2{,}84\times10^4}{50}
\approx5{,}7\times10^2\,\text{W}\!\cdot\!\text{m}^{-2}.
\]

\textbf{4.} Użyteczna moc chłodnicza klimatyzatora wynosi
\[
P_{\text{cool}}=\mathrm{COP}\,P_{\mathrm{el}}
=3{,}0\times2{,}0\times10^3
=6{,}0\times10^3\,\text{W}.
\]
Na jednostkę powierzchni pomieszczenia przypada
\[
\frac{P_{\text{cool}}}{A_{\text{pom}}}
=\frac{6{,}0\times10^3}{20}
=3{,}0\times10^2\,\text{W}\!\cdot\!\text{m}^{-2}.
\]
W tym modelu moc ewapotranspiracji drzewa na jednostkę powierzchni jest prawie dwukrotnie większa niż dla typowego klimatyzatora.

\textbf{5.} Porównanie dotyczy jedynie strumieni energii. Klimatyzator chłodzi zamkniętą, kontrolowaną przestrzeń, natomiast para z drzewa rozprasza się w atmosferze, a wiatr doprowadza ciepłe powietrze. Odczuwalny efekt zależy od wentylacji, wilgotności, temperatury otoczenia i odległości od drzewa. Transpiracja zależy też od nasłonecznienia, wiatru, wilgotności powietrza, gatunku i dostępności wody, a podczas suszy może silnie maleć. Drzewo daje ponadto cień, bezpośrednio zmniejszając promieniowanie słoneczne docierające do gruntu i ludzi. Obliczone moce pozwalają więc porównać rzędy wielkości, ale nie oznaczają jednakowego komfortu cieplnego.

\textbf{6.} W praktyce efekt chłodzący roślin doniczkowych byłby bardzo mały. U większości roślin światło sprzyja otwieraniu aparatów szparkowych, czyli niewielkich porów liści, przez które uchodzi para wodna. Oświetlenie wewnątrz pomieszczeń jest zazwyczaj słabe, więc aparaty szparkowe otwierają się mniej i transpiracja maleje. W słabo wentylowanym pomieszczeniu wzrost wilgotności dodatkowo stopniowo hamuje parowanie.

Wiele roślin może zatem wywołać niewielkie lokalne chłodzenie ewaporacyjne, lecz nie zastąpi klimatyzatora. Aby trwale usuwać energię z mieszkania, trzeba byłoby odprowadzać wilgotne powietrze na zewnątrz. Intensywne oświetlenie sztuczne mogłoby pobudzać transpirację, ale zużyta energia elektryczna niemal w całości zamieniłaby się w ciepło w pomieszczeniu. Wyjątek stanowią przystosowane do suszy rośliny CAM, w tym kaktusy: ich aparaty szparkowe otwierają się głównie nocą.
\end{solution}
\end{exo}

\begin{exo}[Mieszanie dwóch mas wody]{calorimetrie-melange-eau}
\seoready{true}
\keywords{kalorymetria, mieszanie, pojemność cieplna, ciepło właściwe, Joseph Black, temperatura równowagi}

Do naczynia zawierającego $m_2=0{,}300\,\text{kg}$ wody o temperaturze $T_2=15\,^\circ\text{C}$ wlewa się $m_1=0{,}200\,\text{kg}$ wody o temperaturze $T_1=80\,^\circ\text{C}$. Naczynie jest idealnym kalorymetrem: pomija się straty ciepła do otoczenia oraz pojemność cieplną samego naczynia. Przypomnijmy zależność $Q=mc\Delta T$, gdzie $c=4{,}18\times10^3\,\text{J}\!\cdot\!\text{kg}^{-1}\!\cdot\!\text{K}^{-1}$ jest ciepłem właściwym ciekłej wody.

\begin{enumerate}
  \item Uzasadnić, że ciepło oddane przez wodę gorącą zostaje w całości pobrane przez wodę zimną. Zapisać równanie zachowania z wielkościami $m_1$, $m_2$, $c$, $T_1$, $T_2$ i temperaturą równowagi $T_{eq}$.
  \item Wyprowadzić wyrażenie na $T_{eq}$ i obliczyć jego wartość liczbową.
  \item Obliczyć ciepło $Q$ oddane przez gorącą wodę podczas mieszania.
  \item Czy mieszanie dwóch porcji tej samej substancji pozwala wyznaczyć $c$? Uzasadnić.
\end{enumerate}

\begin{indication}
Kalorymetr jest izolowany i sztywny, dlatego całkowita energia wewnętrzna $U_1+U_2$ jest zachowana: ciepło oddane przez jedną porcję zostaje pobrane przez drugą z przeciwnym znakiem.
\end{indication}

\begin{solution}
\textbf{1.} Z zasady zachowania energii wewnętrznej wynika
\[
m_1c(T_{eq}-T_1)+m_2c(T_{eq}-T_2)=0.
\]
\textbf{2.} Wspólny czynnik $c$ skraca się, zatem
\[
T_{eq}=\frac{m_1T_1+m_2T_2}{m_1+m_2}
=\frac{0{,}200\times80+0{,}300\times15}{0{,}200+0{,}300}
=41\,^\circ\text{C}.
\]
\textbf{3.} Gorąca woda oddaje
\[
Q=m_1c(T_1-T_{eq})
=0{,}200\times4{,}18\times10^3\times(80-41)
\approx3{,}26\times10^4\,\text{J}.
\]
Zimna woda pobiera dokładnie tę samą ilość:
\[
m_2c(T_{eq}-T_2)=0{,}300\times4{,}18\times10^3\times26
\approx3{,}26\times10^4\,\text{J}.
\]
\textbf{4.} Nie. Dla dwóch mas tej samej substancji współczynnik $c$ mnoży oba składniki bilansu i skraca się. Temperatura równowagi nie zależy więc od $c$, lecz jest średnią temperatur początkowych ważoną masami. Wyznaczenie ciepła właściwego metodą mieszanin wymaga substancji o znanej pojemności cieplnej, jak w następnym zadaniu.
\end{solution}
\end{exo}

\begin{exo}[Wyznaczanie ciepła właściwego metalu metodą mieszanin]{calorimetrie-methode-melanges}
\seoready{true}
\keywords{kalorymetria, metoda mieszanin, ciepło właściwe, kalorymetr, równoważnik wodny}

Podobnie jak Joseph Black w XVIII wieku, chcemy wyznaczyć ciepło właściwe $c_m$ nieznanego metalu \emph{metodą mieszanin}. Próbkę o masie $m=0{,}150\,\text{kg}$ ogrzewa się do $T_m=95\,^\circ\text{C}$, a następnie szybko zanurza w kalorymetrze zawierającym $m_e=0{,}250\,\text{kg}$ wody o $c_e=4{,}18\times10^3\,\text{J}\!\cdot\!\text{kg}^{-1}\!\cdot\!\text{K}^{-1}$ i temperaturze początkowej $T_e=18\,^\circ\text{C}$. Zmierzona temperatura równowagi wynosi $T_{eq}=22\,^\circ\text{C}$. Początkowo pomija się pojemność cieplną naczynia, mieszadła i termometru.

\begin{enumerate}
  \item Zapisać zasadę zachowania energii dla izolowanego układu metal–woda, pozostawiając $c_m$ jako jedyną niewiadomą.
  \item Wyprowadzić wzór na $c_m$ i obliczyć jego wartość. Jakiemu typowemu metalowi odpowiada ona w przybliżeniu? Wartości odniesienia w $\text{J}\!\cdot\!\text{kg}^{-1}\!\cdot\!\text{K}^{-1}$: aluminium $9{,}0\times10^2$, żelazo $4{,}5\times10^2$, miedź $3{,}9\times10^2$, ołów $1{,}3\times10^2$.
  \item Naczynie kalorymetru pochłania część ciepła. Opisujemy to za pomocą \emph{równoważnika wodnego} $\mu=1{,}5\times10^{-2}\,\text{kg}$: kalorymetr zachowuje się jak dodatkowa masa wody $\mu$. Ponownie obliczyć $c_m$ i omówić kierunek poprawki.
  \item Dlaczego próbkę należy zanurzyć szybko, a kalorymetr dobrze odizolować od otaczającego powietrza?
\end{enumerate}

\begin{indication}
Metal oddaje $Q_m=mc_m(T_m-T_{eq})$, w całości pobrane przez wodę oraz, w pytaniu 3, przez kalorymetr: $Q_m=(m_e+\mu)c_e(T_{eq}-T_e)$.
\end{indication}

\begin{solution}
\textbf{1.} Dla układu izolowanego
\[
mc_m(T_m-T_{eq})=m_ec_e(T_{eq}-T_e).
\]
\textbf{2.} Stąd
\[
c_m=\frac{m_ec_e(T_{eq}-T_e)}{m(T_m-T_{eq})}
=\frac{0{,}250\times4{,}18\times10^3\times(22-18)}{0{,}150\times(95-22)}
\approx3{,}82\times10^2\,\text{J}\!\cdot\!\text{kg}^{-1}\!\cdot\!\text{K}^{-1}.
\]
Wartość ta jest bardzo bliska wartości dla miedzi.

\textbf{3.} Po uwzględnieniu równoważnika wodnego
\[
c_m=\frac{(m_e+\mu)c_e(T_{eq}-T_e)}{m(T_m-T_{eq})}
=\frac{(0{,}250+0{,}015)\times4{,}18\times10^3\times4}{0{,}150\times73}
\approx4{,}05\times10^2\,\text{J}\!\cdot\!\text{kg}^{-1}\!\cdot\!\text{K}^{-1}.
\]
Poprawka zwiększa $c_m$: pominięcie naczynia prowadziło do zaniżenia ciepła rzeczywiście oddanego przez metal, a więc i jego ciepła właściwego.

\textbf{4.} Metoda zakłada układ izolowany przez cały czas pomiaru. Szybkie zanurzenie ogranicza wymianę ciepła gorącej próbki z powietrzem, zanim trafi ona do wody; dobra izolacja ogranicza straty podczas ustalania się równowagi. Każdy nieuwzględniony przepływ ciepła zafałszowałby bilans i wyznaczoną wartość $c_m$. Taką właśnie staranność eksperymentalną stosowali już Black, a później Lavoisier i Laplace w kalorymetrze lodowym.
\end{solution}
\end{exo}

\begin{exo}[Kalorie spożywcze i moc metaboliczna]{calorimetrie-calories-alimentaires}
\seoready{true}
\keywords{kaloria, kilokaloria, Kaloria spożywcza, równoważnik mechaniczny, moc metaboliczna, jednostki}

Etykieta żywieniowa podaje, że dorosły człowiek spożywa średnio około $2000\,\text{kcal}$ dziennie. Kilokaloria jest jednostką spoza SI nadal używaną w dietetyce; jej przeliczenie na SI wynosi $1\,\text{kcal}=4{,}18\times10^3\,\text{J}$.

\begin{enumerate}
  \item Przeliczyć dzienną porcję energii na dżule.
  \item Wyznaczyć średnią moc metaboliczną w watach, zakładając równomierne zużycie energii przez 24 godziny.
  \item Baton energetyczny ma oznaczenie „$210\,\text{kcal}$”. Jaką masę wody o temperaturze początkowej $20\,^\circ\text{C}$ można byłoby doprowadzić do wrzenia w $100\,^\circ\text{C}$ za pomocą tej energii, gdyby została w całości zamieniona na ciepło? Przyjąć $c_{\text{wody}}=4{,}18\times10^3\,\text{J}\!\cdot\!\text{kg}^{-1}\!\cdot\!\text{K}^{-1}$.
  \item W jakich postaciach organizm człowieka rzeczywiście wykorzystuje tę energię? Odpowiedzieć jakościowo.
\end{enumerate}

\begin{indication}
W pytaniu 2 użyć $\mathcal P=E/\Delta t$. W pytaniu 3 najpierw przeliczyć energię na dżule, a następnie wyznaczyć $m$ z $Q=mc\Delta T$.
\end{indication}

\begin{solution}
\textbf{1.} W jednostkach SI
\[
E=2000\times4{,}18\times10^3=8{,}36\times10^6\,\text{J}.
\]
\textbf{2.} Średnia moc wynosi
\[
\mathcal P=\frac{8{,}36\times10^6}{8{,}64\times10^4}\approx97\,\text{W}.
\]
W ciągu całej doby człowiek wykorzystuje więc średnio moc zbliżoną do mocy żarówki. Ten często zaskakujący rząd wielkości pokazuje użyteczność przeliczania kalorii spożywczych na typowe jednostki fizyczne.

\textbf{3.} Energia batona wynosi
\[
Q=210\times4{,}18\times10^3\approx8{,}78\times10^5\,\text{J}.
\]
Dla $\Delta T=80\,\text{K}$,
\[
m=\frac{Q}{c\Delta T}
=\frac{8{,}78\times10^5}{4{,}18\times10^3\times80}
\approx2{,}63\,\text{kg}.
\]
Jeden baton zawiera więc, co do rzędu wielkości, dość energii, aby doprowadzić do wrzenia około $2{,}5\,\text{kg}$ wody z kranu.

\textbf{4.} Organizm nie „zużywa” energii w sensie jej zanikania, lecz przekształca energię chemiczną składników odżywczych. Część służy pracy mięśni, działaniu narządów, utrzymaniu gradientów elektrycznych komórek, transportowi cząsteczek i syntezom chemicznym. Część może zostać zmagazynowana jako glikogen lub tłuszcz. Ostatecznie większość wykorzystanej energii rozprasza się jako ciepło, pomagając utrzymać temperaturę ciała, po czym jest przekazywana otoczeniu. Niewielka część opuszcza też organizm jako energia chemiczna zawarta w produktach przemiany materii.
\end{solution}
\end{exo}
